
\documentclass[12pt,a4paper]{article}

\usepackage[T2A]{fontenc}
\usepackage[utf8]{inputenc}
\usepackage[russian]{babel}
\usepackage{amsmath,amssymb,bm}
\usepackage[a4paper,margin=20mm]{geometry}
\usepackage{mathrsfs}

\newcommand{\vb}[1]{\bm{#1}}
\newcommand{\B}{\vb{B}}
\newcommand{\uvec}{\vb{u}}
\newcommand{\VA}{\vb{V}_A}
\newcommand{\PA}{P_A}
\newcommand{\Rref}{\mathcal R}
\newcommand{\Id}{\mathbf I}

\begin{document}

\section{Обозначения}

Искомые функции времени $t$ и пространственных координат
$\vb{x}$: $\rho$ --- плотность плазмы,
$\uvec$ --- скорость, $p$ --- давление,
$\B$ --- магнитное поле,
$E$ --- полная плотность энергии плазмы и магнитного поля
(без энергии волн), $w_+$ и $w_-$ --- плотности энергии
альфвеновских волн, распространяющихся вдоль
и против $\B$ соответственно.

\[
E=\frac{\rho|\uvec|^2}{2}
+\frac{p}{\gamma-1}
+\frac{|\B|^2}{2\mu_0},
\qquad \gamma=\frac53.
\]
\[
\VA=\frac{\B}{\sqrt{\mu_0\rho}},
\qquad
\PA=\frac{w_++w_-}{2},
\qquad
\vb{g}=-\frac{GM_\odot}{r^3}\vb{x},
\quad r=|\vb{x}|.
\]

Здесь $\VA$ --- вектор скорости Альфвена,
$\PA$ --- волновое давление,
$\mu_0$ --- магнитная постоянная,
$G$ --- гравитационная постоянная,
$M_\odot$ --- масса Солнца.

\section{Трёхмерная система уравнений}

\begingroup
\small
\begin{align}
\frac{\partial\rho}{\partial t}+\nabla\cdot(\rho\uvec)
&=0,
\label{eq:mass}\\[5pt]
\frac{\partial(\rho\uvec)}{\partial t}+\nabla\cdot\left[\rho\uvec\otimes\uvec+\left(p+\PA+\frac{|\B|^2}{2\mu_0}\right)\Id-\frac{\B\otimes\B}{\mu_0}\right]
&=\rho\vb{g},
\label{eq:momentum}\\[5pt]
\frac{\partial\B}{\partial t}+\nabla\cdot(\B\otimes\uvec-\uvec\otimes\B)
&=\vb{0}, \quad \nabla\cdot\B=0
\label{eq:induction}\\[5pt]
\frac{\partial E}{\partial t}+\nabla\cdot\left[\left(E+p+\frac{|\B|^2}{2\mu_0}\right)\uvec-\frac{(\uvec\cdot\B)\B}{\mu_0}+\vb{q}\right]
&=\rho\uvec\cdot\vb{g}-\uvec\cdot\nabla\PA+Q_w-Q_{\rm rad}+Q_{\rm bg},
\label{eq:energy}\\[5pt]
\frac{\partial w_+}{\partial t}+\nabla\cdot\left[(\uvec+\VA)w_+\right]+\frac{w_+}{2}\nabla\cdot\uvec
&=-\Rref\sqrt{w_+w_-}-\Gamma_+w_+,
\label{eq:wplus}\\[5pt]
\frac{\partial w_-}{\partial t}+\nabla\cdot\left[(\uvec-\VA)w_-\right]+\frac{w_-}{2}\nabla\cdot\uvec
&=\Rref\sqrt{w_+w_-}-\Gamma_-w_-.
\label{eq:wminus}
\end{align}
\endgroup

Уравнения \eqref{eq:mass} и \eqref{eq:momentum} описывают
сохранение массы и импульса, \eqref{eq:induction} ---
магнитную индукцию в приближении идеальной МГД,
\eqref{eq:energy} --- баланс энергии плазмы и магнитного поля,
\eqref{eq:wplus}--\eqref{eq:wminus} --- перенос энергии
альфвеновских волн.

Здесь $\Id$ --- единичный тензор,
$\otimes$ --- диадное произведение,
$\vb{q}$ --- тепловой поток,
$Q_w$ --- турбулентный нагрев,
$Q_{\rm rad}$ --- радиационные потери,
$Q_{\rm bg}$ --- дополнительный нагрев,
$\Rref$ --- коэффициент отражения волн,
$\Gamma_\pm$ --- коэффициенты турбулентной диссипации.

\medskip
\noindent\textbf{Турбулентное замыкание.}

\[
\Gamma_\pm=\frac{2}{L_\perp}
\sqrt{\frac{w_\mp}{\rho}},
\qquad
L_\perp\sqrt{|\B|}=C_L=\mathrm{const}>0, \quad Q_w=\Gamma_+w_++\Gamma_-w_-. 
\]

Здесь $L_\perp$ --- поперечный корреляционный масштаб
турбулентности. Нагрев $Q_w$ равен суммарной скорости
диссипации энергии обеих популяций волн.

\medskip
\noindent\textbf{Отражение альфвеновских волн.}

\[
\mathcal R_{\rm imb}
=
\sqrt{
\left[(\VA\cdot\nabla)\ln|\VA|\right]^2
+
\left[\vb{b}\cdot(\nabla\times\uvec)\right]^2
},
\qquad
\vb{b}=\frac{\B}{|\B|}.
\]

\[
\Rref=
\min\left\{
\mathcal R_{\rm imb},
\max(\Gamma_+,\Gamma_-)
\right\}\mathcal C,
\]
\[
\mathcal C=
\begin{cases}
1-2\sqrt{\dfrac{w_-}{w_+}},
& w_+\geq 4w_-,\\[8pt]
0,
& \dfrac14w_-<w_+<4w_-,\\[8pt]
2\sqrt{\dfrac{w_+}{w_-}}-1,
& w_-\geq 4w_+.
\end{cases}
\]

Здесь $\mathcal R_{\rm imb}$ --- максимальная скорость
отражения, определяемая неоднородностью фонового течения,
$\mathcal C$ --- коэффициент, учитывающий дисбаланс
энергий встречных волн. При $\Rref>0$ энергия передаётся от $w_+$ к $w_-$,
при $\Rref<0$ --- от $w_-$ к $w_+$.

Замыкания для теплопроводности $\vb{q}$,
радиационных потерь $Q_{\rm rad}$,
дополнительного нагрева $Q_{\rm bg}$ будут рассмотрены позднее.


\section{Квазиодномерная система уравнений}
\label{sec:quasi1d}

Рассматривается течение плазмы вдоль магнитной трубки
с площадью поперечного сечения $A(r)$.
Боковая поверхность трубки $\Sigma$ образована магнитными
силовыми линиями и является непроницаемой для плазмы:
\[
\left.\B\cdot\vb{n}\right|_\Sigma=0,
\qquad
\left.\uvec\cdot\vb{n}\right|_\Sigma=0,
\qquad
\uvec\parallel\B.
\]
Здесь $\vb{n}$ --- единичная нормаль к боковой поверхности.
Предполагается, что все искомые величины зависят только
от продольной координаты $r$ и времени $t$,
а магнитная трубка имеет заданную стационарную геометрию.

\begingroup
\small
\begin{align}
\frac{\partial\rho}{\partial t}
+\frac{1}{A}\frac{\partial(A\rho u)}{\partial r}
&=0,
\label{eq:1d_mass}\\[5pt]
\frac{\partial(\rho u)}{\partial t}
+\frac{1}{A}\frac{\partial}{\partial r}
\left[A(\rho u^2+p+\PA)\right]
&=\frac{p+\PA}{A}\frac{dA}{dr}
-\rho\frac{GM_\odot}{r^2},
\label{eq:1d_momentum}\\[5pt]
\frac{\partial H}{\partial t}
+\frac{1}{A}\frac{\partial}{\partial r}
\left[A\bigl(u(H+p)+q\bigr)\right]
&=-\rho u\frac{GM_\odot}{r^2}
-u\frac{\partial\PA}{\partial r}
+Q_w-Q_{\rm rad}+Q_{\rm bg},
\label{eq:1d_energy}\\[5pt]
\frac{\partial w_+}{\partial t}
+\frac{1}{A}\frac{\partial}{\partial r}
\left[A(u+V_A)w_+\right]
+\frac{w_+}{2A}\frac{\partial(Au)}{\partial r}
&=-\Rref\sqrt{w_+w_-}-\Gamma_+w_+,
\label{eq:1d_wplus}\\[5pt]
\frac{\partial w_-}{\partial t}
+\frac{1}{A}\frac{\partial}{\partial r}
\left[A(u-V_A)w_-\right]
+\frac{w_-}{2A}\frac{\partial(Au)}{\partial r}
&=\Rref\sqrt{w_+w_-}-\Gamma_-w_-.
\label{eq:1d_wminus}
\end{align}
\endgroup

Уравнения \eqref{eq:1d_mass}--\eqref{eq:1d_wminus}
описывают соответственно сохранение массы,
продольного импульса, энергии плазмы и перенос
энергии альфвеновских волн.

Плотность энергии плазмы и волновое давление:
\[
H=\frac{\rho u^2}{2}+\frac{p}{\gamma-1},
\qquad
\PA=\frac{w_++w_-}{2},
\qquad \gamma=\frac53.
\]

\medskip
\noindent\textbf{Геометрия и магнитное поле.}

Из сохранения магнитного потока через поперечное
сечение трубки следует
\[
A(r)B(r)=\Phi_B=\mathrm{const},
\qquad
V_A=\frac{B(r)}{\sqrt{\mu_0\rho}}.
\]

Здесь $A(r)$ --- площадь поперечного сечения
магнитной трубки, $\Phi_B$ --- постоянный магнитный поток,
$B(r)$ --- модуль магнитного поля, $V_A$ ---
скорость Альфвена. Конкретная зависимость $A(r)$
будет задана отдельно.

\medskip
\noindent\textbf{Турбулентная диссипация.}

\[
\Gamma_\pm=\frac{2}{L_\perp}
\sqrt{\frac{w_\mp}{\rho}},
\qquad
L_\perp\sqrt{B}=C_L=\mathrm{const}>0, \quad Q_w=\Gamma_+w_++\Gamma_-w_-.
\]

Здесь $L_\perp$ --- поперечный корреляционный масштаб
турбулентности, $\Gamma_\pm$ --- коэффициенты
диссипации энергии волн, $Q_w$ --- нагрев плазмы
за счёт турбулентной диссипации.

\medskip
\noindent\textbf{Отражение альфвеновских волн.}

В квазиодномерном приближении коэффициент отражения
принимает вид
\[
\Rref=
\min\left\{
\left|\frac{\partial V_A}{\partial r}\right|,
\max(\Gamma_+,\Gamma_-)
\right\}\mathcal C,
\]
где $\mathcal C$ --- коэффициент дисбаланса энергий
встречных волн, определённый в трёхмерной модели.

\medskip
\noindent\textbf{Остальные источники энергии.}

В уравнении \eqref{eq:1d_energy}
$q$ --- продольный тепловой поток,
$Q_{\rm rad}$ --- радиационные потери,
$Q_{\rm bg}$ --- дополнительный нагрев плазмы.

Замыкания для $q$, $Q_{\rm rad}$ и $Q_{\rm bg}$
будут заданы отдельно.


\section{Упрощённая модель с альфвеновским нагревом}
\label{sec:simplified}

На первом этапе рассмотрим упрощённую модель,
в которой пренебрегаем теплопроводностью,
радиационными потерями и дополнительным нагревом:
\[
q=0,\qquad Q_{\rm rad}=0,\qquad Q_{\rm bg}=0.
\]
Единственным механизмом нагрева плазмы считается
турбулентная диссипация альфвеновских волн.
При этом сохраняются перенос и отражение волн,
а также их взаимодействие со средним течением.

Система уравнений принимает вид:

\begingroup
\small
\begin{align}
\frac{\partial\rho}{\partial t}
+\frac{1}{A}\frac{\partial(A\rho u)}{\partial r}
&=0,
\label{eq:s_mass}\\[5pt]
\frac{\partial(\rho u)}{\partial t}
+\frac{1}{A}\frac{\partial}{\partial r}
\left[A(\rho u^2+p+\PA)\right]
&=\frac{p+\PA}{A}\frac{dA}{dr}-\rho\frac{GM_\odot}{r^2},
\label{eq:s_momentum}\\[5pt]
\frac{\partial H}{\partial t}
+\frac{1}{A}\frac{\partial}{\partial r}
\left[Au(H+p)\right]
&=-\rho u\frac{GM_\odot}{r^2}
-u\frac{\partial\PA}{\partial r}+Q_w,
\label{eq:s_energy}\\[5pt]
\frac{\partial w_+}{\partial t}
+\frac{1}{A}\frac{\partial}{\partial r}
\left[A(u+V_A)w_+\right]
+\frac{w_+}{2A}\frac{\partial(Au)}{\partial r}
&=-\Rref\sqrt{w_+w_-}-\Gamma_+w_+,
\label{eq:s_wplus}\\[5pt]
\frac{\partial w_-}{\partial t}
+\frac{1}{A}\frac{\partial}{\partial r}
\left[A(u-V_A)w_-\right]
+\frac{w_-}{2A}\frac{\partial(Au)}{\partial r}
&=\Rref\sqrt{w_+w_-}-\Gamma_-w_-.
\label{eq:s_wminus}
\end{align}
\endgroup

Все обозначения и замыкания для $\Gamma_\pm$,
$\Rref$ и $Q_w$ остаются прежними.
В дальнейшем исследуем возможность получения
стационарного решения этой системы и его использование
в качестве начального условия для нестационарной задачи.


\section{Стационарная система уравнений}
\label{sec:stationary}

Рассмотрим стационарное течение, полагая все производные
по времени равными нулю. Система уравнений принимает вид:

\begingroup
\small
\begin{align}
\frac{d}{dr}(A\rho u)
&=0,
\label{eq:st_mass}\\[5pt]
\frac{1}{A}\frac{d}{dr}
\left[A(\rho u^2+p+\PA)\right]
&=\frac{p+\PA}{A}\frac{dA}{dr}-\rho\frac{GM_\odot}{r^2},
\label{eq:st_momentum}\\[5pt]
\frac{1}{A}\frac{d}{dr}\left[Au(H+p)\right]
&=-\rho u\frac{GM_\odot}{r^2}
-u\frac{d\PA}{dr}+Q_w,
\label{eq:st_energy}\\[5pt]
\frac{1}{A}\frac{d}{dr}\left[A(u+V_A)w_+\right]
+\frac{w_+}{2A}\frac{d(Au)}{dr}
&=-\Rref\sqrt{w_+w_-}-\Gamma_+w_+,
\label{eq:st_wplus}\\[5pt]
\frac{1}{A}\frac{d}{dr}\left[A(u-V_A)w_-\right]
+\frac{w_-}{2A}\frac{d(Au)}{dr}
&=\Rref\sqrt{w_+w_-}-\Gamma_-w_-.
\label{eq:st_wminus}
\end{align}
\endgroup

\noindent\textbf{Замыкания.}

\[
H=\frac{\rho u^2}{2}+\frac{p}{\gamma-1},
\qquad
\PA=\frac{w_++w_-}{2},
\qquad \gamma=\frac53.
\]

\[
AB=\Phi_B=\mathrm{const},
\qquad
V_A=\frac{B}{\sqrt{\mu_0\rho}},
\qquad
L_\perp\sqrt B=C_L=\mathrm{const}>0.
\]

\[
\Gamma_\pm=\frac{2}{L_\perp}\sqrt{\frac{w_\mp}{\rho}},
\qquad
Q_w=\Gamma_+w_++\Gamma_-w_-.
\]

\[
\Rref=
\min\left\{
\left|\frac{dV_A}{dr}\right|,
\max(\Gamma_+,\Gamma_-)
\right\}\mathcal C.
\]

Коэффициент $\mathcal C$ определяется ранее заданной
кусочной функцией от $w_+$ и $w_-$.
Геометрия магнитной трубки задаётся функцией $A(r)$.

Из уравнения неразрывности следует первый интеграл
\[
A\rho u=\dot m=\mathrm{const},
\]
где $\dot m$ --- постоянный массовый расход через трубку.

\newpage

\section{Обезразмеривание стационарной задачи}

Рассмотрим стационарную систему для пяти функций
\[
\rho(r),\quad u(r),\quad p(r),\quad w_+(r),\quad w_-(r).
\]
Температура связана с давлением и плотностью
уравнением состояния идеального газа с постоянной $\mu$:
\[
p=\frac{\rho k_{\rm B}T}{\mu m_u},\qquad
H=\frac{\rho u^2}{2}+\frac{p}{\gamma-1},
\qquad \gamma=\frac53.
\]

Параметры задачи и заданная геометрия записываются как
\[
G,\ M_\odot,\ \mu_0,\ k_{\rm B},\ m_u,\ \mu,\ \gamma,\
R_\odot,\ r_{\rm out},\
\rho_{\rm ph},\ T_{\rm ph},\ B_{\rm ph},\
L_{\perp,{\rm ph}},\ w_{+,{\rm ph}}
\qquad\text{и}\qquad A(r).
\]
При $B_{\rm ph}>0$ магнитный поток и корреляционная
длина определяются соотношениями
\[
AB=A(R_\odot)B_{\rm ph},\qquad
L_\perp\sqrt B=L_{\perp,{\rm ph}}\sqrt{B_{\rm ph}}.
\]

\medskip
\noindent\textbf{Характерные величины.}

Используем следующие основные размерности СИ
\[
\mathsf M\ \text{(масса)},\quad
\mathsf L\ \text{(длина)},\quad
\mathsf T\ \text{(время)},\quad
\mathsf I\ \text{(сила тока)},\quad
\Theta\ \text{(температура)}.
\]
Выбираем масштабы
\[
\begin{array}{c|c|c}
\text{Величина} & \text{Масштаб} & \text{Размерность}\\ \hline
\text{Длина} & R_*=R_\odot & \mathsf L\\
\text{Скорость} & U_*=\sqrt{GM_\odot/R_*}
    & \mathsf L\mathsf T^{-1}\\
\text{Плотность} & \rho_*=\rho_{\rm ph}
    & \mathsf M\mathsf L^{-3}\\
\text{Температура} & T_*=T_{\rm ph}
    & \Theta\\
\text{Площадь} & A_*=A(R_*)
    & \mathsf L^2\\
\text{Давление и плотность энергии}
    & p_*=\rho_*U_*^2
    & \mathsf M\mathsf L^{-1}\mathsf T^{-2}\\
\text{Магнитное поле} & B_*=B_{\rm ph}
    & \mathsf M\mathsf T^{-2}\mathsf I^{-1}\\
\text{Частота} & \Gamma_*=U_*/R_*
    & \mathsf T^{-1}\\
\text{Объёмный нагрев} & Q_*=p_*U_*/R_*
    & \mathsf M\mathsf L^{-1}\mathsf T^{-3}
\end{array}
\]
Гравитационный масштаб скорости обращает коэффициент
гравитации в единицу. Фотосферные масштабы плотности
и температуры упрощают граничные условия, а общий
масштаб давления и энергий волн сохраняет вид
соотношения для волнового давления.

Безразмерные величины определим следующим образом:
\[
\hat r=\frac r{R_*},\quad
\hat A=\frac A{A_*},\quad
\hat\rho=\frac\rho{\rho_*},\quad
\hat u=\frac u{U_*},\quad
\hat T=\frac T{T_*},\quad
\hat B=\frac B{B_*},\quad
\hat V_A=\frac{V_A}{U_*},
\]
\[
\hat p=\frac p{p_*},\quad
\hat H=\frac H{p_*},\quad
\hat w_\pm=\frac{w_\pm}{p_*},\quad
\hat P_A=\frac{P_A}{p_*},\quad
\hat L_\perp=\frac{L_\perp}{R_*},
\]
\[
\hat\Gamma_\pm=\frac{\Gamma_\pm}{\Gamma_*},\qquad
\hat{\mathcal R}=\frac{\mathcal R}{\Gamma_*},\qquad
\hat Q_w=\frac{Q_w}{Q_*}.
\]

\medskip
\noindent\textbf{Подстановка масштабов.}

В члене силы тяжести возникает множитель
$GM_\odot/(R_*U_*^2)=1$.
Подстановка масштабов в уравнение состояния,
скорость Альфвена и частоты диссипации даёт
\[
\hat p
=\frac{k_{\rm B}T_*}{\mu m_uU_*^2}\hat\rho\hat T,
\qquad
\hat V_A
=\frac{B_*}{U_*\sqrt{\mu_0\rho_*}}\,
\frac{1}{\hat A\sqrt{\hat\rho}},
\qquad
\hat\Gamma_\pm
=\frac{2R_*}{L_{\perp,{\rm ph}}}
\sqrt{\frac{\hat w_\mp}{\hat A\hat\rho}}.
\]
Обозначим получившиеся коэффициенты через
\[
\theta=\frac{k_{\rm B}T_*}{\mu m_uU_*^2},\qquad
b=\frac{B_*}{U_*\sqrt{\mu_0\rho_*}},\qquad
d=\frac{2R_*}{L_{\perp,{\rm ph}}}.
\]

Далее шапочки опускаем. Все переменные в уравнениях
и граничных условиях ниже безразмерны, а параметры
с индексом ${\rm ph}$ в определениях масштабов
и коэффициентов остаются размерными.

\medskip
\noindent\textbf{Безразмерная система.}

Штрих обозначает производную по безразмерному $r$.
Стационарные уравнения имеют вид
\begingroup
\small
\begin{align*}
(A\rho u)'&=0,\\[4pt]
\frac1A\left[A(\rho u^2+p+P_A)\right]'
&=(p+P_A)\frac{A'}A-\frac{\rho}{r^2},\\[4pt]
\frac1A\left[Au(H+p)\right]'
&=-\frac{\rho u}{r^2}-uP_A'+Q_w,\\[4pt]
\frac1A\left[A(u+V_A)w_+\right]'
+\frac{w_+}{2A}(Au)'
&=-\mathcal R\sqrt{w_+w_-}-\Gamma_+w_+,\\[4pt]
\frac1A\left[A(u-V_A)w_-\right]'
+\frac{w_-}{2A}(Au)'
&=\mathcal R\sqrt{w_+w_-}-\Gamma_-w_-.
\end{align*}
\endgroup
Здесь
\[
H=\frac{\rho u^2}{2}+\frac{p}{\gamma-1},
\qquad
P_A=\frac{w_++w_-}{2},
\qquad
V_A=\frac{b}{A\sqrt\rho},
\]
\[
\Gamma_\pm=d\sqrt{\frac{w_\mp}{A\rho}},
\qquad
Q_w=\Gamma_+w_++\Gamma_-w_-
=\frac{d}{\sqrt{A\rho}}
\left(w_+\sqrt{w_-}+w_-\sqrt{w_+}\right).
\]
Коэффициент отражения принимает вид
\[
\mathcal R
=\mathcal C(w_+,w_-)
\min\left\{
b\left|\left(\frac{1}{A\sqrt\rho}\right)'\right|,
\frac{d}{\sqrt{A\rho}}
\max\bigl(\sqrt{w_+},\sqrt{w_-}\bigr)
\right\}.
\]
Вид функции $\mathcal C$ сохраняется, поскольку
она зависит от отношения энергий волн.

Температуру восстанавливаем по найденным
давлению и плотности:
\[
T=\frac{p}{\theta\rho}.
\]

\medskip
\noindent\textbf{Физический смысл коэффициентов.}

Коэффициент
\[
\theta
=\frac{k_{\rm B}T_{\rm ph}/(\mu m_u)}{GM_\odot/R_\odot}
=\frac{p_{\rm ph}}{\rho_{\rm ph}U_*^2}
\]
сравнивает фотосферный тепловой масштаб $p_{\rm ph}/\rho_{\rm ph}$
с гравитационным масштабом удельной энергии.
При записи системы через $p$ он входит в граничное
значение давления и в формулу пересчёта температуры.

Коэффициент
\[
b=\frac{V_{A,{\rm ph}}}{U_*}
\]
есть отношение фотосферной скорости Альфвена
к гравитационному масштабу скорости.
Он входит в перенос волн и в градиентную ветвь
коэффициента отражения.

Коэффициент
\[
d=\frac{2R_\odot}{L_{\perp,{\rm ph}}}
\]
есть удвоенное отношение радиуса Солнца
к фотосферной корреляционной длине.
Он входит в диссипацию волн и нагрев газа.
Безразмерная частота $\Gamma_\pm$ представляет
отношение времени переноса $R_*/U_*$ к местному
времени диссипации соответствующей популяции.

\medskip
\noindent\textbf{Граничные условия.}

На фотосфере задаются размерные плотность
$\rho_{\rm ph}$, температура $T_{\rm ph}$
и энергия выходящих волн $w_{+,{\rm ph}}$.
Граничное давление вычисляем по уравнению состояния:
\[
p_{\rm ph}=\frac{\rho_{\rm ph}k_{\rm B}T_{\rm ph}}{\mu m_u}.
\]
Для системы, записанной через давление, это даёт
\[
\rho(1)=1,\qquad
p(1)=\theta,\qquad
w_+(1)=\omega,\qquad
\omega=\frac{w_{+,{\rm ph}}}{\rho_{\rm ph}U_*^2}.
\]
При этих условиях $T(1)=1$.
Параметр $\omega$ сравнивает заданную волновую
энергию с масштабом давления $\rho_{\rm ph}U_*^2$.
Отношение $\omega/(2\theta)$ показывает вклад
выходящих волн в давление относительно теплового
давления на фотосфере.

Значения $u(1)$ и $w_-(1)$ подбираются при поиске
стационарного решения, регулярно проходящего через
звуковую критическую точку. Для определения обоих
значений нужны две независимые скалярные связи.
Регулярность в простой звуковой критической точке
даёт одну связь, а возможность получить вторую
из другой критической точки проверяется
при анализе связанной системы.

Внешняя граница находится при $r=r_{\rm out}/R_\odot$.
На ней граничные данные не задаются, если все
характеристики выходят из расчётной области.
Для волновой подсистемы проверяется также $u>V_A$.

\section{Особая точка и построение стационарного решения}

Все величины безразмерны, $\gamma=5/3$.
Геометрия $A(r)$ задана.

\subsection{Сведение системы}

Постоянные массового расхода и удельного полного
потока энергии определяются соотношениями
\[
m=A\rho u,
\]
\[
\varepsilon
=\frac{u^2}{2}-\frac1r+\frac{5p}{2\rho}
+\frac{3(w_++w_-)}{2\rho}
+\frac{V_A(w_+-w_-)}{\rho u}.
\]
Следовательно,
\[
\rho=\frac{m}{Au},\qquad
V_A=b\sqrt{\frac{u}{mA}},\qquad
\Gamma_\pm=d\sqrt{\frac{uw_\mp}{m}},
\]
\begin{equation}
p=\frac{2m}{5Au}
\left(\varepsilon-\frac{u^2}{2}+\frac1r\right)
-\frac35(w_++w_-)
-\frac{2V_A}{5u}(w_+-w_-).
\label{eq:crit_pressure}
\end{equation}

Таким образом, остаются три искомые функции
$u,w_+,w_-$ и постоянные $m,\varepsilon$.
Давление далее рассматривается как функция
$r,u,w_+,w_-$, заданная формулой
\eqref{eq:crit_pressure}.

\subsection{Три ветви отражения}

Замыкание отражения имеет вид
\[
\mathcal R
=\mathcal C\min\left\{
\frac{V_A}{2}\left|\frac{u'}u-\frac{A'}A\right|,
\max(\Gamma_+,\Gamma_-)
\right\}.
\]
На каждой ветви оно линейно по $u'$:
\[
\mathcal R=\eta u'+\zeta.
\]
Соответствующие коэффициенты равны
\[
\begin{array}{c|c|c}
\text{Ветвь} & \eta & \zeta\\ \hline
\text{Ограничение диссипацией}
&0&\mathcal C\max(\Gamma_+,\Gamma_-)\\[3pt]
\text{Положительный градиент}
&\dfrac{\mathcal C V_A}{2u}
&-\dfrac{\mathcal C V_A}{2}\dfrac{A'}A\\[7pt]
\text{Отрицательный градиент}
&-\dfrac{\mathcal C V_A}{2u}
&\dfrac{\mathcal C V_A}{2}\dfrac{A'}A
\end{array}
\]

Для первой ветви требуется
\[
\frac{V_A}{2}\left|\frac{u'}u-\frac{A'}A\right|
\geq\max(\Gamma_+,\Gamma_-).
\]
Для второй и третьей ветвей требуется обратное
неравенство и соответственно
\[
\frac{u'}u-\frac{A'}A\geq0,
\qquad
\frac{u'}u-\frac{A'}A\leq0.
\]
Эти условия проверяются после нахождения $u'$.

\subsection{Уравнение для производной скорости}

После исключения плотности и давления система
принимает вид
\begin{equation}
\mathsf M
\begin{pmatrix}u'\\w_+'\\w_-'\end{pmatrix}
=\boldsymbol f,
\label{eq:crit_matrix}
\end{equation}
где
\[
\mathsf M=
\begin{pmatrix}
\rho u+\dfrac{\partial p}{\partial u}
&
\dfrac{\partial p}{\partial w_+}+\dfrac12
&
\dfrac{\partial p}{\partial w_-}+\dfrac12\\[8pt]
\dfrac{w_+}{2}\left(3+\dfrac{V_A}{u}\right)
+\eta\sqrt{w_+w_-}
&u+V_A&0\\[8pt]
\dfrac{w_-}{2}\left(3-\dfrac{V_A}{u}\right)
-\eta\sqrt{w_+w_-}
&0&u-V_A
\end{pmatrix},
\]
\[
\boldsymbol f=
\begin{pmatrix}
-\rho/r^2-\partial p/\partial r\\[4pt]
-\dfrac{w_+}{2}(3u+V_A)\dfrac{A'}A
-\zeta\sqrt{w_+w_-}-\Gamma_+w_+\\[7pt]
-\dfrac{w_-}{2}(3u-V_A)\dfrac{A'}A
+\zeta\sqrt{w_+w_-}-\Gamma_-w_-
\end{pmatrix}.
\]
Частные производные давления вычисляются по
\eqref{eq:crit_pressure} при фиксированных
$m,\varepsilon$. Производная по $r$ учитывает $A(r)$.

Пусть $u\ne V_A$.
Из второй и третьей строк непосредственно выражаем
\begin{equation}
w_+'=\frac{f_2-M_{21}u'}{u+V_A},
\qquad
w_-'=\frac{f_3-M_{31}u'}{u-V_A}.
\label{eq:crit_wave_derivatives}
\end{equation}
Подставляя их в первую строку, получаем
\begin{equation}
D\,u'=N,
\label{eq:crit_scalar}
\end{equation}
где
\[
D=M_{11}
-\frac{M_{12}M_{21}}{u+V_A}
-\frac{M_{13}M_{31}}{u-V_A},
\]
\[
N=f_1
-\frac{M_{12}f_2}{u+V_A}
-\frac{M_{13}f_3}{u-V_A}.
\]
Таким образом, в обычной точке $u'=N/D$.
Если $D=0$, конечная производная возможна
только при одновременном выполнении $N=0$.

\subsection{Поиск критического состояния}

Условия особой точки имеют вид
\begin{equation}
D(r_c,u_c,w_{+,c},w_{-,c};m,\varepsilon)=0,
\qquad
N(r_c,u_c,w_{+,c},w_{-,c};m,\varepsilon)=0.
\label{eq:crit_conditions}
\end{equation}

При заданных $A(r),b,d$ здесь присутствуют
шесть неизвестных
\[
r_c,\quad u_c,\quad w_{+,c},\quad w_{-,c},
\quad m,\quad\varepsilon.
\]
Поэтому значения $m,\varepsilon$ сами по себе
обычно не определяют критическую точку.

Для локального поиска можно задать также
$w_{+,c},w_{-,c}$.
Тогда \eqref{eq:crit_conditions} становится системой
двух алгебраических уравнений для $r_c,u_c$.

Другой способ --- сначала построить кривые
$D=0$ и $N=0$ на плоскости $(r,u)$.
Их пересечения дают начальные приближения
для уточнения корней.
Пересечений может быть несколько или не быть вовсе.

Локальный поиск проводится отдельно для трёх
ветвей отражения.
Каждая ветвь даёт свои функции $D,N$.
Для каждого найденного состояния далее вычисляется
наклон и проверяются неравенства соответствующей
ветви. Недопустимые кандидаты отбрасываются.
Совпадающие решения на границах ветвей учитываются
как одна точка.

Если $u=V_A$, выражения с делением на $u-V_A$
не используются. В этом случае проверяются
вырождение и совместность исходной системы
\eqref{eq:crit_matrix}.

\subsection{Наклон в особой точке}

Обозначим искомый наклон скорости через $q=u'_c$.
Тогда
\[
w_{+,c}'=\frac{f_2-M_{21}q}{u_c+V_{A,c}},
\qquad
w_{-,c}'=\frac{f_3-M_{31}q}{u_c-V_{A,c}}.
\]
Все коэффициенты берутся в найденном
критическом состоянии.

Продифференцируем \eqref{eq:crit_scalar}
вдоль решения:
\[
D'u'+Du''=N'.
\]
В особой точке $D=N=0$, поэтому
\begin{equation}
q\left(
D_r+D_uq+D_{w_+}w_{+,c}'+D_{w_-}w_{-,c}'
\right)
=
N_r+N_uq+N_{w_+}w_{+,c}'+N_{w_-}w_{-,c}'.
\label{eq:crit_slope}
\end{equation}
Индексы здесь обозначают частные производные
при фиксированных $m,\varepsilon$.

Поскольку $w_{\pm,c}'$ линейны по $q$,
уравнение \eqref{eq:crit_slope} имеет степень
не выше второй относительно $q$.
Каждый действительный корень задаёт кандидат
на наклон решения.

Для каждого корня проверяются условия выбранной
ветви отражения и направление продолжения решения.
Для простого пересечения требуется также
\[
D_r+D_uq+D_{w_+}w_{+,c}'+D_{w_-}w_{-,c}'\ne0.
\]
Если допустимых действительных корней нет,
выбранное критическое состояние не даёт
гладкого проходящего решения.
При вырождении уравнения для наклона необходимо
разложение более высокого порядка.

На границах ветвей замыкания производные
рассматриваются односторонне.

\subsection{Интегрирование от найденной точки}

После выбора наклона начинаем расчёт
с небольшим отступом от особой точки:
\[
u(r_c\pm\delta)=u_c\pm\delta u_c'+O(\delta^2),
\]
\[
w_\pm(r_c\pm\delta)
=w_{\pm,c}\pm\delta w_{\pm,c}'+O(\delta^2).
\]
От $r_c-\delta$ интегрируем к фотосфере,
от $r_c+\delta$ --- к внешней границе.
В регулярных точках производные находим
из линейной системы \eqref{eq:crit_matrix}.

Ветвь отражения определяется локально
и может изменяться вдоль решения.
При её смене пересчитываются $\eta,\zeta$
и матрица системы.
Проверяются положительность давления и плотности,
а также устойчивость результата при уменьшении
стартового отступа $\delta$.


\end{document}
